-[7] Lớp 7

----[K] Khoa học tự nhiên

-------[0] Mở đầu

----------[1] Phương pháp và kĩ năng học tập môn Khoa học tự nhiên
-------------[1] Phân biệt thời tiết và khí hậu
-------------[2] Các yếu tố khí hậu (nhiệt độ, độ ẩm, lượng mưa, gió)
-------------[3] Đọc hiểu bản tin dự báo thời tiết, ứng dụng vào sinh hoạt hàng ngày
-------------[4] Tìm hiểu đặc điểm khí hậu địa phương và ảnh hưởng đến sản xuất nông nghiệp
-------------[5] Các loại thiên tai thường gặp ở Việt Nam (bão, lũ, hạn hán, động đất, sạt lở...)
-------------[6] Nguyên nhân và hậu quả của thiên tai
-------------[7] Biện pháp phòng chống và ứng phó thiên tai
-------------[8] Xây dựng kế hoạch ứng phó thiên tai cho gia đình và trường học
-------------[9] Biến đổi khí hậu: nguyên nhân, hậu quả và hành động của học sinh

-------[1] Nguyên tử. Sơ lược về bảng tuần hoàn các nguyên tố hoá học

----------[1] Nguyên tử
-------------[1] Nêu mô hình nguyên tử (vỏ, hạt nhân)
-------------[2] Xác định các loại hạt (p, n, e) trong nguyên tử
-------------[3] So sánh khối lượng của hạt nhân với khối lượng của electron
-------------[4] Vận dụng kiến thức về nguyên tử giải thích tại sao nguyên tử trung hoà về điện

----------[2] Nguyên tố hoá học
-------------[1] Phân biệt nguyên tử và nguyên tố hoá học
-------------[2] Viết kí hiệu hoá học của nguyên tố
-------------[3] Tên gọi, kí hiệu của 20 nguyên tố đầu tiên
-------------[4] Bài tập về số hiệu nguyên tử và số khối

----------[3] Sơ lược về bảng tuần hoàn các nguyên tố hoá học
-------------[1] Mô tả cấu tạo của bảng tuần hoàn (ô, chu kì, nhóm)
-------------[2] Xác định vị trí (ô, chu kì, nhóm) của một nguyên tố trong bảng tuần hoàn (với 20 nguyên tố đầu)
-------------[3] Dựa vào vị trí suy ra một số thông tin cơ bản của nguyên tố
-------------[4] Phân biệt kim loại, phi kim, khí hiếm dựa vào vị trí và tính chất
-------------[5] Bìa toán tìm 2 nguyên tố trong cùng chu kì liên tiếp, hai chu kì liên tiếp nhau trong cùng nhóm A
-------------[6] Hoàn thành bảng thông tin về tên nguyên tố, kí hiệu hóa học, số e, p,n, khối lượng nguyên tử

-------[2] Phân tử. Liên kết hoá học

----------[1] Phân tử – Đơn chất – Hợp chất
-------------[1] Lý thuyết - Định nghĩa, khái niệm
-------------[2] Phân loại - Nhận dạng đơn chất/hợp chất
-------------[3] Công thức - Viết và đọc CTHH
-------------[4] Định lượng cơ bản - Tính khối lượng phân tử
-------------[5] Phân tích - Thành phần, tỉ lệ nguyên tố
-------------[6] Tổng hợp - Vận dụng, liên hệ thực tế

----------[2] Giới thiệu về liên kết hoá học
-------------[1] Lý thuyết - Định nghĩa khái niệm liên kết
-------------[2] Cấu hình electron bền vững - Quy tắc octet
-------------[3] Nhận dạng - Phân loại loại liên kết
-------------[4] Sơ đồ hình thành - Mô tả quá trình (định tính) sự hình thành phân tử (sự kết hợp của các nguyên tử)
-------------[5] Sơ đồ hình thành - Mô tả quá trình (định tính) hình thành liên kết ion (cho và nhận electron),cộng hoá trị (dùng chung electron)
-------------[6] Tính toán electron - Số electron nhường/nhận/dùng chung (định lượng)
-------------[7] Tính chất - So sánh chất ion và cộng hóa trị
-------------[8] Tổng hợp - Vận dụng, liên hệ thực tế

----------[3] Hoá trị và công thức hoá học
-------------[1] Lý thuyết - Định nghĩa hóa trị, quy ước
-------------[2] Đọc công thức - Phân tích thông tin từ CTHH
-------------[3] Tìm hóa trị - Từ công thức xác định hóa trị (định tính)
-------------[4] Lập công thức - Từ hóa trị lập CTHH (định lượng cơ bản)
-------------[5] Tổng hợp - Kết hợp tìm hóa trị và lập công thức
-------------[6] Nâng cao - Tính \% khối lượng nguyên tố (định lượng)
-------------[7] Xác định công thức hoá học dựa vào phần trăm khối lượng nguyên tố
-------------[8] Bài tập vận dụng quy tắc hoá trị trong thực tế

-------[3] Tốc độ

----------[1] Tốc độ chuyển động
-------------[1] Khái niệm tốc độ, công thức tính tốc độ (v = s/t)
-------------[2] Đơn vị đo tốc độ (m/s, km/h) và chuyển đổi đơn vị
-------------[3] Dụng cụ đo tốc độ (tốc kế, đồng hồ bấm giây)
-------------[4] Tính tốc độ trong sinh hoạt (đi bộ, xe đạp, xe máy, ô tô...)
-------------[5] Đọc hiểu biển báo tốc độ giao thông và ý nghĩa an toàn
-------------[6] Tính thời gian di chuyển từ nhà đến trường, lập kế hoạch đi du lịch

----------[2] Đo tốc độ

----------[3] Đồ thị quãng đường – thời gian
-------------[1] Cách vẽ và đọc đồ thị quãng đường - thời gian
-------------[2] Xác định tốc độ từ đồ thị
-------------[3] Phân tích hành trình thực tế qua đồ thị (xe buýt, tàu hỏa, chuyến đi gia đình...)

----------[4] Thảo luận về ảnh hưởng của tốc độ trong an toàn giao thông

-------[4] Âm thanh

----------[1] Sóng âm
-------------[1] Nguồn âm và sự truyền âm
-------------[2] Môi trường truyền âm (rắn, lỏng, khí, chân không)
-------------[3] Biên độ, tần số, độ to, độ cao của âm
-------------[4] Giải thích hiện tượng âm thanh trong đời sống (tiếng vang, sấm sét nghe sau chớp...)
-------------[5] Ứng dụng truyền âm: ống nghe y tế, siêu âm, sonar tàu ngầm

----------[2] Độ to và độ cao của âm

----------[3] Phản xạ âm, chống ô nhiễm tiếng ồn
-------------[1] Khái niệm phản xạ âm, vật phản xạ âm tốt/hấp thụ âm tốt
-------------[2] Ô nhiễm tiếng ồn: nguồn gây ô nhiễm và tác hại
-------------[3] Biện pháp chống ô nhiễm tiếng ồn
-------------[4] Thiết kế giải pháp giảm tiếng ồn tại nhà/trường (cách âm phòng học, trồng cây xanh...)
-------------[5] Đánh giá mức độ ô nhiễm tiếng ồn tại khu vực sinh sống (gần chợ, đường quốc lộ, công trường...)

-------[5] Ánh sáng

----------[1] Năng lượng ánh sáng. Tia sáng, vùng tối
-------------[1] Hiện tượng khúc xạ ánh sáng
-------------[2] Thấu kính hội tụ và thấu kính phân kì
-------------[3] Giải thích tại sao đũa cắm trong nước trông như bị gãy
-------------[4] Ứng dụng thấu kính: kính lúp, kính cận, kính viễn, máy ảnh
-------------[5] Liên hệ cấu tạo mắt người với máy ảnh, giải thích cận thị và viễn thị
-------------[6] Ánh sáng trắng và sự phân tích ánh sáng trắng
-------------[7] Sự trộn các ánh sáng màu
-------------[8] Giải thích hiện tượng cầu vồng sau mưa
-------------[9] Ứng dụng ánh sáng màu: đèn LED, màn hình TV, trang trí, đèn giao thông

----------[2] Sự phản xạ ánh sáng
-------------[1] Sự truyền thẳng của ánh sáng
-------------[2] Sự phản xạ ánh sáng, định luật phản xạ
-------------[3] Ảnh của vật qua gương phẳng
-------------[4] Giải thích hiện tượng nhìn thấy vật, bóng tối, nhật thực, nguyệt thực
-------------[5] Ứng dụng gương phẳng trong đời sống (gương soi, kính chiếu hậu, gương nha khoa...)

----------[3] Ảnh của vật qua gương phẳng

-------[6] Từ

----------[1] Nam châm
-------------[1] Tính chất của nam châm (hút sắt, định hướng Bắc-Nam)
-------------[2] Các loại nam châm (nam châm tự nhiên, nam châm nhân tạo)
-------------[3] Ứng dụng nam châm trong đời sống (la bàn, loa, động cơ điện, nam châm tủ lạnh...)
-------------[4] Thực hành sử dụng la bàn xác định phương hướng khi đi dã ngoại

----------[2] Từ trường
-------------[1] Khái niệm từ trường, đường sức từ
-------------[2] Từ trường của nam châm và dòng điện (nam châm điện)
-------------[3] Ứng dụng nam châm điện: chuông điện, cần cẩu điện từ, khóa cửa từ
-------------[4] Giải thích hoạt động của một số thiết bị sử dụng từ trường (máy phát điện, loa, micro...)

----------[3] Chế tạo nam châm điện đơn giản

-------[7] Trao đổi chất và chuyển hoá năng lượng ở sinh vật

----------[1] Khái quát về trao đổi chất và chuyển hoá năng lượng
-------------[1] Khái niệm trao đổi chất ở sinh vật
-------------[2] Vai trò của trao đổi chất đối với cơ thể
-------------[3] Liên hệ thực tế: tại sao cần ăn uống đầy đủ, cung cấp nước cho cơ thể

----------[2] Quang hợp ở thực vật
-------------[1] Khái niệm quang hợp, phương trình quang hợp
-------------[2] Các yếu tố ảnh hưởng đến quang hợp (ánh sáng, nước, CO₂)
-------------[3] Vai trò của quang hợp đối với tự nhiên
-------------[4] Ứng dụng kiến thức quang hợp trong trồng trọt (bố trí khoảng cách cây, chiếu sáng nhà kính...)
-------------[5] Giải thích tại sao nên trồng nhiều cây xanh, tại sao lá cây có màu xanh

----------[3] Một số yếu tố ảnh hưởng đến quang hợp

----------[4] Thực hành: Chứng minh quang hợp ở cây xanh

----------[5] Hô hấp tế bào
-------------[1] Khái niệm hô hấp tế bào, phương trình hô hấp
-------------[2] Mối quan hệ giữa quang hợp và hô hấp
-------------[3] Giải thích tại sao khi tập thể dục nặng thở gấp, tại sao không nên để nhiều cây trong phòng ngủ đóng kín

----------[6] Một số yếu tố ảnh hưởng đến hô hấp tế bào

----------[7] Thực hành: Hô hấp ở thực vật

----------[8] Trao đổi khí ở sinh vật
-------------[1] Trao đổi khí ở thực vật (khí khổng)
-------------[2] Trao đổi khí ở động vật và người (phổi, mang...)
-------------[3] Ứng dụng: tại sao cần thông thoáng phòng học, tác hại khói thuốc lá đến phổi

----------[9] Vai trò của nước và chất dinh dưỡng đối với sinh vật

----------[A] Trao đổi nước và chất dinh dưỡng ở thực vật
-------------[1] Sự hấp thụ nước và muối khoáng ở rễ
-------------[2] Sự vận chuyển nước và chất dinh dưỡng trong thân
-------------[3] Ứng dụng trong trồng trọt: tưới nước hợp lý, bón phân đúng cách, thủy canh

----------[B] Trao đổi nước và chất dinh dưỡng ở động vật
-------------[1] Dinh dưỡng ở động vật (tiêu hóa, hấp thụ)
-------------[2] Vai trò của nước và chất dinh dưỡng đối với cơ thể người
-------------[3] Xây dựng chế độ ăn uống khoa học, cân đối dinh dưỡng cho học sinh
-------------[4] Phân tích thành phần dinh dưỡng trên bao bì thực phẩm

----------[C] Thực hành: Chứng minh thân vận chuyển nước và lá thoát hơi nước

-------[8] Cảm ứng ở sinh vật

----------[1] Cảm ứng ở sinh vật và tập tính ở động vật
-------------[1] Khái niệm cảm ứng, các hình thức cảm ứng ở thực vật
-------------[2] Tính hướng sáng, hướng nước, hướng trọng lực
-------------[3] Quan sát và giải thích hiện tượng cây nghiêng về phía có ánh sáng, rễ mọc hướng xuống
-------------[4] Ứng dụng hiểu biết về cảm ứng thực vật trong trồng cây (xoay chậu cây, bố trí cây trồng...)
-------------[5] Khái niệm cảm ứng ở động vật
-------------[6] Các giác quan ở người (thị giác, thính giác, xúc giác, khứu giác, vị giác)
-------------[7] Phản xạ có điều kiện và không điều kiện
-------------[8] Liên hệ thực tế: phản xạ rụt tay khi chạm nóng, chảy nước miếng khi thấy thức ăn ngon
-------------[9] Bảo vệ các giác quan trong đời sống (bảo vệ mắt, tai khi dùng thiết bị điện tử...)

----------[2] Vận dụng hiện tượng cảm ứng ở sinh vật vào thực tiễn

----------[3] Thực hành: Cảm ứng ở sinh vật

-------[9] Sinh trưởng và phát triển ở sinh vật

----------[1] Khái quát về sinh trưởng và phát triển ở sinh vật
-------------[1] Khái niệm sinh trưởng và phát triển
-------------[2] Các giai đoạn sinh trưởng, phát triển ở thực vật (nảy mầm, ra hoa, kết quả...)
-------------[3] Các yếu tố ảnh hưởng đến sinh trưởng (nước, ánh sáng, nhiệt độ, chất dinh dưỡng)
-------------[4] Ứng dụng trong trồng trọt: chọn thời vụ gieo trồng, sử dụng chất kích thích sinh trưởng
-------------[5] Thực hành theo dõi sự nảy mầm và sinh trưởng của cây đậu tại nhà
-------------[6] Các giai đoạn sinh trưởng, phát triển ở động vật (biến thái hoàn toàn, không hoàn toàn)
-------------[7] Các giai đoạn phát triển ở người (sơ sinh, thiếu niên, dậy thì, trưởng thành...)
-------------[8] Các yếu tố ảnh hưởng đến sinh trưởng ở người (dinh dưỡng, vận động, giấc ngủ)
-------------[9] Chế độ dinh dưỡng và luyện tập hợp lý cho học sinh tuổi dậy thì
-------------[A] Giải thích vòng đời côn trùng ứng dụng trong phòng trừ sâu bệnh (muỗi, bướm...)

----------[2] Ứng dụng sinh trưởng và phát triển ở sinh vật vào thực tiễn

----------[3] Thực hành: Quan sát, mô tả sự sinh trưởng và phát triển ở một số sinh vật

-------[A] Sinh sản ở sinh vật

----------[1] Sinh sản vô tính ở sinh vật
-------------[1] Sinh sản vô tính ở thực vật (giâm, chiết, ghép cành, nuôi cấy mô)
-------------[2] Sinh sản hữu tính ở thực vật (thụ phấn, thụ tinh, tạo quả và hạt)
-------------[3] Ứng dụng nhân giống vô tính trong nông nghiệp (chiết cành bưởi, ghép cây hoa hồng...)
-------------[4] Vai trò của côn trùng trong thụ phấn, bảo vệ ong mật

----------[2] Sinh sản hữu tính ở sinh vật
-------------[1] Sinh sản vô tính ở động vật (nảy chồi, phân mảnh...)
-------------[2] Sinh sản hữu tính ở động vật (thụ tinh trong, thụ tinh ngoài)
-------------[3] Khái quát về sinh sản ở người (cơ quan sinh dục, sự thụ tinh, phát triển phôi thai)
-------------[4] Giáo dục sức khỏe sinh sản tuổi vị thành niên (vệ sinh, phòng tránh xâm hại...)

----------[3] Một số yếu tố ảnh hưởng và điều hoà, điều khiển sinh sản ở sinh vật

----------[4] Cơ thể sinh vật là một thể thống nhất
